\documentclass[10pt,a4paper]{amsart}
\usepackage[margin=19mm]{geometry}
\usepackage{amsmath,amssymb,mathtools}
\usepackage[scr=rsfs,cal=euler]{mathalfa}
\usepackage{xfrac}
\usepackage[unicode,hidelinks,bookmarks=false]{hyperref}
\newcommand{\ie}{i.e.\ }
\newcommand{\cf}{cf.\ }
\newcommand{\resp}{respectively\ }
\newcommand{\Subsection}{\subsection}
\providecommand{\nobreakdash}{\nobreak\hbox{-}}
\setlength{\emergencystretch}{2em}
\multlinegap0pt
\usepackage{bm}
\usepackage{bbm}
\usepackage{dsfont}
\usepackage{xspace}
\usepackage{cancel}
\usepackage{mathtools}
\usepackage[shortlabels]{enumitem}
\usepackage{amsthm}
\makeatletter
\@ifundefined{proofclaim}{\newtheorem{proofclaim}{Proofclaim}}{}
\@ifundefined{theorem}{\newtheorem{theorem}{Theorem}}{}
\@ifundefined{claim}{\newtheorem{claim}[theorem]{Claim}}{}
\makeatother
\newcommand{\id}{\mathcal{I}}
\newcommand{\li}{\mathbb{L}(\mathcal{I}^+)}
\newcommand{\nwd}{\mathsf{nwd}}
\newcommand{\restr}[2]{#1\upharpoonright {#2}}
\title[Properties of Laver forcing associated with a co-ideal expre]{Properties of Laver forcing associated with a co-ideal expressed via the Katetov order}
\author{Osvaldo Guzm\'an, Francisco Santiago Nieto-de la Rosa, Ulises Ariet Ramos-Garcia}
\date{Source: arXiv:2601.09061v2 (CC BY 4.0)}
\begin{document}
\maketitle

\noindent\emph{Source and licence.} Properties of Laver forcing associated with a co-ideal expressed via the Katetov order. Version arXiv:2601.09061v2 (CC BY 4.0), distributed under the Creative Commons Attribution 4.0 International licence, as recorded in the arXiv licence metadata for this exact version and shown on the abstract page https://arxiv.org/abs/2601.09061v2. Full rights notice: licenses/arxiv-2601.09061-CC-BY-4.0.txt.\par\smallskip

\noindent\textit{Editorial scope.} Counted unit: the Theorem whose source label is \texttt{Cohen reals} in the article (source0.tex lines 218--220, source label \texttt{Cohen reals}), delivered together with the article's own complete original proof of it (source0.tex lines 221--331, one \texttt{proof} environment of 111 lines). No mathematics is written, repaired or re-derived: statement and proof are the article's own text line for line. Required context retained uncounted: none. Every definition, display and label of those lines is retained unchanged. Omitted from this package and not redistributed: 1--217, 332--1370. Nothing inside a retained excerpt is omitted or altered: every retained body is delivered exactly as the article's lines are written, so no omission receipt of any kind accompanies this package. Citations: the retained excerpts carry no citation command at all, so no bibliography entry of the article is transcribed and no reference is redistributed. Reference adaptation: none was needed and none was made; every reference inside the delivered unit resolves inside it, so no unresolved reference or citation is produced. Printed slips kept literal: no printed word, symbol, hypothesis or number of the counted statement or of its proof was altered. Layout and class: the article's own class is \texttt{amsart} with packages including mathscinet, amssymb, latexsym, stmaryrd, mathtools, enumerate, mathrsfs, hyperref, cmll, xcolor, tikz; the wrapper is the lane's pinned \texttt{amsart} editorial wrapper at 10pt on A4, which restates the theorem-like environments the retained text needs and adds the article's own macro definitions that the retained text needs (\texttt{\textbackslash id}, \texttt{\textbackslash li}, \texttt{\textbackslash nwd}, \texttt{\textbackslash restr}), with extra wrapper packages (bm, bbm, dsfont, xspace, cancel, mathtools, enumitem). The delivered numbering is the unit's own. Assets: no retained span contains a figure environment, a graphics-inclusion command, a PDF-page inclusion, a table or a TikZ picture; no image file is copied into this package, the package declares no asset dependency and no image file is redistributed with it. Licence: version arXiv:2601.09061v2 of this article is distributed under the Creative Commons Attribution 4.0 International licence, as recorded in the arXiv licence metadata for this exact version and shown on the abstract page https://arxiv.org/abs/2601.09061v2. A full rights notice is delivered at licenses/arxiv-2601.09061-CC-BY-4.0.txt. Characters: every retained excerpt is pure ASCII, so the delivered engine, XeLaTeX with its default Latin Modern text font, reports no missing character.
\begin{theorem} \label{Cohen reals}
    For every ideal $\id$ on $\omega$. Then $\li$ adds Cohen reals if and only if there is $X\in\id^+$ such that $\nwd \leq_K \restr{\id}{X}$.
\end{theorem}

\begin{proof}
$(\Leftarrow)$ To prove that $\li$ adds a Cohen real, fix $X \in \id^+$ and a witness $f: X \to 2^\omega$ for $\nwd \leq_K \id \upharpoonright X$. Set $P = X^{<\omega}$. Recursively we ca construct a $\li$-name $\dot c$ such that for every $T \leq P$, if $k$ is the maximum such that $T$ decides $\dot c \upharpoonright k$, then for all $n\in succ_T(stem(T))$,

\begin{displaymath}
    T_{stem(T)^\frown n}\Vdash `` \forall i<n \, ( \dot c(k+i)=f(n)(i))".
\end{displaymath}
    
We claim that $\dot c$ is a Cohen real.

Let $N \in \nwd$, and let $R_N$ be the nowhere dense tree such that $N = [R_N]$. For any $T \leq P$, we show there exists $T' \leq T$ with $T' \Vdash ``\dot c \notin N"$. Let $k < \omega$ be the largest number such that $T$ decides $\dot c \upharpoonright k$, and let $s \in 2^k$ with $T \Vdash ``\dot c \upharpoonright k = s"$.

For each $n \in succ_T(stem(T))$, define $g_n : \omega \to 2$ such that for all $i<k$, $g_n(i)=s(i)$, and for $i \geq k$, $g_n(i)=f(n)(i-k)$. Notice that $g_n$ is a translation of $f(n)$, since $f[succ_T(stem(T))]\in \nwd^+$, then $G=\{g_n:n\in succ_T(stem(T))\}\in \nwd^+$. Since $N\in \nwd$ and $int(\bar{G})\not=\emptyset$ there is $s'\supseteq s$ such that $\langle s'\rangle\subseteq \bar G$ and $s'\notin R_N$. In particular $\langle s'\rangle\cap G\not=\emptyset$. In addition, if \[Y=\{n\in succ_T(stem(T)):s'\subseteq g_n\},\] then $Y$ is infinite.

Fix $n\in Y$ larger than $k$ and $|s'|$. By construction
\begin{displaymath}
    T_{stem(T)^\frown n}\Vdash``\forall i<n (\dot c(k+i)=f(n)(i))".
\end{displaymath}
Choosing $i=|s'|$ we have
\begin{displaymath}
    T_{stem(T)^\frown n}\Vdash``\dot c(k+i)=f(n)(i)=g_n(i+k)".
\end{displaymath}

Since $s'\notin R_N$ we have that $T_{stem(T)^\frown n}$ forces $\dot c\upharpoonright k+i+1\not\in R_N$ and we are done.

$(\Rightarrow)$ By contrapositive. Let $\dot h$ be a $\li$- name for a real and $T\in \li$. We must show that there is $S\leq T$ and $A$ a nowhere dense tree such that $S\Vdash``\dot h\in [A]"$.

%%%%
%For each $k$ let $L_{k+1}=\{s_k^\frown n:n<\omega\}$.%%%%

For all $R\in \li$ enumerate $R\setminus \{\sigma\in R:\sigma\leq stem(R)\}$ as $\{\sigma_k^R:k\in \omega\}$ such that:

\begin{enumerate}
    \item[(a)] $\sigma_k^R<\sigma_\ell^R$ implies $k<\ell$ and
    \item[(b)] if $\sigma$ is such that $\sigma_k^R=\sigma^\frown n$, $\sigma_\ell^R=\sigma^\frown m$ and $n<m$, then $k<\ell$.
\end{enumerate}

In order to obtain the condition $S$ and the nowhere dense set $A$, we construct by recursion for each $k<\omega$ the following:

\begin{enumerate}[(i)]
    \item $\sigma_k^S\in\omega^{<\omega}$.
    \item $S^k\in \li$.
    \item $X_{k}\in \id^+$.
    \item $\forall n\in X_k, \{T^{\ell}(k,n):\ell<\omega\} \subseteq \li$.
    \item $f_k:X_{k}\rightarrow2^\omega$.
    \item $B_{k}$ a nowhere dense tree.
    \item $g(k),n_k, \ell_k\in\omega$ .
\end{enumerate}


such that $\{\sigma^S_{k}:k<\omega\}$ has the properties (a) and (b) and for all $k$:

\begin{enumerate}
    \item $\ell_k\leq k$ and $\ell_k=k$ only if $k=0$.
    \item $\sigma^S_0=stem(T)$ and $\sigma^S_{k+1}=\sigma^{S^k}_{k+1}={\sigma^{S^k}_{\ell_k}}^\frown n_k$.
    \item $S^0=\bigcup\{T^{0}(0,n):n\in X_0\}$, $S^{k+1}\leq S^{k}$ and  \\$S^{k+1}=\left(S^k\setminus S^k_{\sigma^S_{k+1}}\right)\cup \left\{T^{g(k)}(\ell_j,n_j):j\in X_{k+1}\right\}$.
    \item $X_0\subseteq succ_T(\sigma_0^S)$ and $X_{k+1}\subseteq succ_{T^{g(k)}(\ell_k,n_k)}(\sigma^S_{k+1})$.
    \item For all $\ell<\omega$ and $n\in X_k$, $T^{\ell}(k,n)$ decides $\dot h \upharpoonright \ell+1$ and
    \\$T^{\ell+1}(k,n)\leq_0 T^\ell(k,n)\leq_0 S^k_{{\sigma_k^S}^\frown n}$ 
    \item $f_k(n)(\ell)=i$ if and only if $T^\ell(k,n)\Vdash``\dot h(\ell)=i"$.
    \item $ f_k[X_k]\subseteq [B_k]$.
    \item $g(k+1)>g(k)$.
    \item For all $\bar h\in f_{k+1}[X_{k+1}]$, $\bar h\upharpoonright g(k)=f_{\ell_k}(n_k)\upharpoonright g(k)$.
    \item For all $k'>k, \,B_{k'}\cap 2^{\leq g(k)}\subseteq \bigcup_{i\leq k+1} B_i$.
    \item For all $t\in 2^{\leq k}$, there is $t'\supseteq t$ such that $t'\in 2^{\leq g(k)}$ and $ t'\notin \bigcup_{i\leq k}B_i$.
\end{enumerate}

Let us make some remarks: $S=\{\sigma: \sigma\subset  stem(T)\}\cup \{\sigma^S_{k}:k<\omega\}$ is, in fact, a condition because $succ_S(\sigma^S_k)=X_k\in\id^+$ for $k<\omega$. The succession mention in (iv) with property (5) can be obtained using the pure decision property. By conditions (3) and (5) we see that $S^{k+1}_{\sigma^{S}_{k+1}}$ decides $\dot h\upharpoonright g(k)$. Conditions (1) and (2) are given by properties (a) and (b). By (2), for all $k$ we have $n_k\in succ_{S^k}(\sigma^{S^k}_{\ell_k})$. We do not specify $n_k$ explicitly; it is determined by (2) and the enumeration of $S^k$.


%Given $k$, if we define $f'_k:succ_{T^{g(k)}(\ell_k,n_k)}(\sigma^S_{k+1})$ with the condition (6), since $f'_k$

%Once we have defined $\{B_i:i<k\}$

%\begin{displaymath}
%    f_0(n)(k)=i \text{ if and only if } T^k(0,n)\Vdash``\dot h(k)=i".
%\end{displaymath}

We start the recursion. Basic step $k=0$. We define $f_0':succ_T(stem(T))\rightarrow2^\omega$ as in (6).
Since $f_0'$ is not a Kat\v{e}tov function, there is $X_0\in \id^+\upharpoonright succ_T(stem(T))$ such that $f_0'[X_0]\subseteq [B_0]$ for some $B_0$ a nowhere dense tree. Take $f_0=f_0'\upharpoonright X_0$. Choose $t'$ such that $t'\not \in B_0$ and define $g(0)=|t'|$.

Recursion step. Assume that our construction has been done for all $i\leq k$. Now do it for $k+1$. Define $f'_{k+1}:succ_{T^{g(k)}(\ell_k,n_k)}(\sigma^S_{k+1}) \rightarrow2^\omega$ as in (6). Since $f_{k+1}'$ is not a Kat\v{e}tov function there is $X_{k+1}\in \id^+$ and $B_{k+1}'$ a nowhere dense tree such that if we define $f_{k+1}=f'_{k+1}\upharpoonright X_{k+1}$, then $f_{k+1}[X_{k+1}]\subseteq B_{k+1}$. We can prune $B_{k+1}'$ to $B_{k+1}$ such that if $t\in B_{k+1}$ has size at most $g(k)$, then there is $n\in X_{k+1}$ such that $f_{k+1}(n)\upharpoonright |t|= t$. Choose $g(k+1)$ in such a way that (11) is satisfied.

Now we need to prove that conditions (9) and (10) hold. To prove (9) take $\bar h\in f_{k+1}[X_{k+1}]$ and $j<g(k)$. Let $m\in X_{k+1}$ be such that $\bar h=f_{k+1}(m)$. By (6), $\bar h(j)=i$ if and only if $T^j(k+1,m)\Vdash``\dot h(j)=i"$. In the same way, $f_{\ell_k}(n_k)(j)=i$ if and only if $T^\ell(\ell_k,n_k)\Vdash``\dot h(j)=i"$.
Notice that by (2),(3) and (5)
\[T^j(k+1,m)\leq S^{k}_{{\sigma^S_{k+1}}^\frown m}=S^{k}_{{\sigma^S_{\ell_k}}^\frown n_k^\frown m}
\]
and  $m\in X_{k+1}\subseteq succ_{T^{g(k)}(\ell_k,n_k)}({\sigma^{S^k}_{\ell_k}}^\frown n_k)$, then we have 
\begin{displaymath}
    T^j(k+1,m)\leq T^{g(k)}(\ell_k,n_k)\leq_0 T^j(\ell_k,n_k).
\end{displaymath}
To prove (10), by the recursion hypothesis we just verify for $k'=k+1$ and $j\leq k$. Let $t\in B_{k+1}\cap2^{\leq g(j)}$. Take $t'\in B_{k+1}\cap2^{g(k)}$ with $t \subseteq t'$. By (9) $t'=f_{\ell_k}(n_k)\upharpoonright g(k)$. Thus $t=f_{\ell_k}(n_k)\upharpoonright |t|$. If $j\geq \ell_k$, $t\in B_{\ell_k}$ and we are done. If $j<\ell_k$ we know that $t\in B_{\ell_k}\cap 2^{g(j)}$. By inductive hypothesis $t\in \bigcup_{i\leq j} B_i$ and we are done.



\begin{claim}
$A=\bigcup\{B_k:k<\omega\}$ is a nowhere dense tree.
\end{claim}
\begin{proofclaim}
    Let $t\in A$, there is $k$ such that $|t|=k$. By (11) there is $t'\in 2^{<g(k)}\setminus (\bigcup_{\ell\leq k}B_{\ell})$. If there is some $k'>k$ with $t'\in B_{k'}$. By (10) there is $\ell\leq k$ with $t'\in B_\ell$. This is an absurd, then $t'\notin B_{k'}$ for all $k'>k$.
\end{proofclaim}

The proof ends when we show that $S\Vdash``\dot h \in [A]$". Searching for a contradiction assume that it does not hold. There are $j<\omega$ and $S'\leq S$ such that $S'\Vdash``\dot h\upharpoonright j\notin A"$. Let $k<\omega$ be such that $stem(S')=\sigma^{S}_k$. Let $k'>k$ be  such that $j<g(k')$ and $\sigma^S_{k'}\in S'$. Then, by (6) and (7),
\begin{displaymath}
    T^j(\ell_{k'}, n_{k'})\Vdash ``\dot h\upharpoonright j\in B_{k'}".
\end{displaymath}
and by(3), 
\[
S'_{\sigma^S_{k'}}\leq S_{\sigma^S_{k'}}\leq S^{k'+1}_{\sigma^S_{k'}}=S^{k'+1}_{{\sigma^S_{\ell_{k'}}}^\frown n_{k'}}\leq_0 T^{g(k')}(\ell_{k'},n_{k'})\leq_0  T^j(\ell_{k'}, n_{k'}).
\]
Thus, $S'_{\sigma^S_{k'}}\Vdash`` \dot h\upharpoonright j\in A"$ which is a contradiction.
\end{proof}

\end{document}
